\subsection{Proposed Global Scheduling Algorithm}
\label{sec:method-ours}

We focus on the widely used setting where, at any timestep $t$,
the sampler ultimately maintains \emph{a single latent and a single noise}
$(x_t, \varepsilon_t)$, but is allowed to perform \emph{multiple local search
iterations} inside this timestep before committing to $x_{t-1}$.
Concretely, at step $t$ we run a local noise search operator $\mathcal{L}_t$
around the current latent $x_t$, and only the best candidate noise is kept and
used to produce $x_{t-1}$ via $F_\theta$.
Under this restriction, the global scheduling problem reduces to a sequence of
\emph{per-timestep decisions}: at each $t$, given the history of search so far
and the remaining budget, should we continue local search at this timestep,
or stop and move on to $t-1$?

This raises an inherent difficulty.
When making the decision at timestep $t$, we do not know:
(1) the potential improvement that additional search at $t$ could still bring,
and (2) the potential improvement that could be obtained by saving NFE and
spending it on later timesteps.
The scheduling strategy below tackles these two aspects
from complementary angles.

\paragraph{Offline coarse time scheduling.}
First, for each model and dataset, we perform an \emph{offline} analysis on a
small set of prompts or images.
We run noise search with a simple local operator $\mathcal{L}_t$ and a uniform,
relatively large number of iterations at every timestep, and log the verifier
scores after each local search iteration.
From these logs we can compute, for each timestep $t$, statistics such as the
average score gain brought by local search and its variability across examples.
Aggregating these statistics over many runs reveals which parts of the
denoising trajectory are consistently more sensitive to noise refinement and
thus more ``valuable'' to search.
In our offline profiling, we observe a clear model-dependent pattern:
for Stable Diffusion (SD), most of the gains concentrate on early timesteps,
whereas for EDM the improvements are concentrated in a contiguous middle segment
of the denoising trajectory.

Based on this analysis, we partition timesteps into a small number of groups,
such as a \emph{high-value} region and a \emph{low-value} region, and assign a
larger share of the total NFE budget to the high-value region.
Within each region, we derive a coarse per-step budget $\{K_t\}$ (e.g., a
higher average number of local iterations in high-value steps and a lower
average in low-value steps).
This yields a time schedule tailored to the model and dataset, fixed across
prompts, that captures \emph{where search tends to matter the most}.

\paragraph{Online fine-grained scheduling.}
Second, at test time, we add an \emph{online} controller that refines how the
pre-assigned budget $\{K_t\}$ is actually used, on a per-prompt basis.
For each timestep $t$, we run local search up to $K_t$ iterations, but allow the
controller to stop earlier or, within a small window, adaptively modulate how
aggressively to spend the budget.

The controller uses two simple statistics computed from the local search
feedback at timestep $t$:
(i) the \emph{score gain} of local search, and
(ii) the \emph{score variance} across candidates within this step.
Let $s_t^{(j)}$ be the best verifier score at step $t$ after $j$ local search
iterations, and let $\{s_{t,\text{cand}}^{(j,i)}\}_i$ denote the scores of all
candidate noises evaluated in iteration $j$.
We define the incremental gain
\[
  g_t^{(j)} = s_t^{(j)} - s_t^{(j-1)}
\]
and the per-iteration variance
\[
  \operatorname{Var}_t^{(j)} = \operatorname{Var}_i\bigl(s_{t,\text{cand}}^{(j,i)}\bigr).
\]
$g_t^{(j)}$ measures whether we are still finding
better noises at step $t$, while $\operatorname{Var}_t^{(j)}$ reflects how
sensitive the verifier is to noise perturbations at this timestep.

Our online policy is based on the following intuition:
if the recent gains $g_t^{(j)}$ remain positive and the score variance
$\operatorname{Var}_t^{(j)}$ is large, then this timestep is both important
and not yet saturated, so we continue local search until we reach the allocated
budget $K_t$ (or a small upper margin).
Conversely, if gains have diminished (e.g., recent $g_t^{(j)}$ are close to
zero or negative) and the variance across candidates is small, then local
search at this timestep is likely exhausted and further NFE is better reserved
for subsequent timesteps; in this case we stop early and move on to $t-1$.
These decisions are made online based solely on statistics observed up to step
$t$, without access to future timesteps. To ensure the realized compute exactly matches the prescribed total NFE budget,
we reconcile any over-/under-spending caused by adaptive stopping/slack:
any extra NFE consumed in the high-value region is compensated by allocating
additional iterations to the first few low-value steps, while any NFE deficit
is corrected by trimming iterations from the last low-value steps.

\paragraph{Overall procedure.}
A single sampling run draws an initial noise $x_T \sim \mathcal{N}(0, I)$ and obtains the
coarse per-step budget $\{K_t\}$ from the offline time schedule.
Then, for $t = T, \dots, 1$, we run the local noise search operator
$\mathcal{L}_t$ around the current latent $x_t$, while the online controller
monitors the sequence of gains $\{g_t^{(j)}\}$ and variances
$\{\operatorname{Var}_t^{(j)}\}$ to decide whether to continue or stop local
search at this timestep, subject to the budget $K_t$.
The final chosen noise at step $t$ is then fed into the base sampler
$F_\theta$ to obtain $x_{t-1}$, and the process repeats until we reach $x_0$.

Under the same total NFE budget, this two-level scheduling concentrates
computation on timesteps that are both globally and instance-wise valuable,
avoiding wasted search on steps where noise perturbations have little
impact on the verifier.

\begin{algorithm}[t]
\caption{GAINS: Global Adaptive Inference-time Noise Scheduling}
\label{alg:offline_online_budget}
\begin{algorithmic}[1]
\STATE \textbf{Input:} timesteps $\{t_T,\dots,t_1\}$, total budget $B$, offline allocation $\{K_t\}$ ($\sum_t K_t=B$),
slack $\delta$, gain window $W_g$, threshold coefs $(\beta_g,\beta_\sigma)$, black-box $\textsc{LocalIter}(t)$.
\STATE \textbf{Output:} realized per-step iterations $\{\widehat K_t\}$.

\STATE $R \leftarrow B$; \quad $\mathcal{H}_g\leftarrow\emptyset$ (historical gains); \quad $\mathcal{H}_\sigma\leftarrow\emptyset$ (historical variances).
\FOR{$t=T,T\!-\!1,\dots,1$}
    \STATE $K \leftarrow \min(K_t, R)$; \quad $K_{\max}\leftarrow \min(K+\delta, R)$; \quad $j_{\text{watch}}\leftarrow \max(1, K-\delta)$.
    \STATE $s^{(0)} \leftarrow -\infty$; \quad $\mathcal{H}\leftarrow\emptyset$ (recent gains); \quad $\widehat K_t \leftarrow 0$.
    \FOR{$j=1,2,\dots,K_{\max}$}
        \STATE $(s^{(j)}, \sigma^{(j)}) \leftarrow \textsc{LocalIter}(t)$
        \STATE $g^{(j)} \leftarrow s^{(j)} - s^{(j-1)}$; push $g^{(j)}$ into $\mathcal{H}$ and keep last $W_g$.
        \STATE $\widetilde g^{(j)} \leftarrow \mathrm{Mean}(\mathcal{H})$.
        \STATE $\widehat K_t \leftarrow \widehat K_t + 1$; \quad $R \leftarrow R - 1$.
        \IF{$R=0$} \STATE \textbf{break} \ENDIF
        \IF{$j \ge j_{\text{watch}}$ \AND $|\mathcal{H}_g|>0$ \AND $|\mathcal{H}_\sigma|>0$}
            \STATE $g_{\mathrm{th}}\leftarrow \beta_g \cdot \mathrm{Mean}(\mathcal{H}_g)$;\quad
                   $\sigma_{\mathrm{th}}\leftarrow \beta_\sigma \cdot \mathrm{Mean}(\mathcal{H}_\sigma)$.
            \IF{$\widetilde g^{(j)} < g_{\mathrm{th}}$ \AND $\sigma^{(j)} < \sigma_{\mathrm{th}}$}
                \STATE \textbf{break}
            \ENDIF
        \ENDIF
    \ENDFOR
    \STATE append $\mathrm{Mean}(\mathcal{H})$ to $\mathcal{H}_g$; append $\sigma^{(\widehat K_t)}$ to $\mathcal{H}_\sigma$.
\ENDFOR
\STATE \textbf{return} $\{\widehat K_t\}$.
\end{algorithmic}
\end{algorithm}
